% GENERATED FILE. Edit canonical lessons, then run npm run generate:books.
% canonical-source-hash: fnv1a64:4335189fbac34936
% canonical-lessons: ES-C667-apartamento, ES-C667-propietario, ES-C667-anfitrion, ES-C667-inquilino, ES-C667-alquilar

\chapter{Un piso de vacaciones}
\label{ch:es-un-piso-de-vacaciones}
\hlchaptermodality{667}

\begin{chapteropening}
\textbf{By the end of this chapter:} \emph{I can rent a holiday flat or a car: meet the owner or host and talk about tenants.}

Complete the last lesson of chapter 667: I can rent a holiday flat or a car: meet the owner or host and talk about tenants.
\end{chapteropening}

\section[el apartamento]{el apartamento — the flat, the apartment}

\label{lesson:ES-C667-apartamento}



\begin{quote}
Before the new one: say the Spanish for the caravan, then the Spanish for the torch.
\end{quote}

\subsection*{You'll want to know: el apartamento}
\textbf{el apartamento} — \textquotedblleft{}the flat, the apartment\textquotedblright{}.

\begin{quote}
Alquilamos un apartamento cerca de la playa.
\end{quote}

\begin{quote}
\emph{We rented a flat near the beach.}
\end{quote}

Spain also says \emph{el piso}; much of Latin America says \emph{el departamento}.

\begin{cousinweb}
\emph{apartamento} is usually traced to Italian \emph{appartamento}, perhaps through French \emph{appartement}; the Italian is built on \emph{appartare} \textquotedblleft{}to set apart\textquotedblright{}, from Latin \emph{ad} and \emph{pars} \textquotedblleft{}part\textquotedblright{} (EtymDB).
\end{cousinweb}

\begin{grammarlens}[title={What you've built}]
One more word for telling, explaining and describing at B1.
\end{grammarlens}

\subsection*{Your turn}
\begin{itemize}
\raggedright
  \item \emph{Say it:} \emph{el apartamento}
  \item \emph{Say it:} \emph{el apartamento}, once more
  \item \emph{Say it:} a sentence of your own with \emph{el apartamento}
  \item \emph{Recall:} say the Spanish for the caravan, then the Spanish for the torch, then say \emph{el apartamento} again
\end{itemize}

\begin{tcolorbox}[breakable,colback=teal!4,colframe=teal!35!black,title={Before you move on}]
What does el apartamento mean? (\textquotedblleft{}The flat, the apartment\textquotedblright{}.) Say it once more.
\end{tcolorbox}
\section[el propietario]{el propietario — the proprietor, the owner; the landlord}

\label{lesson:ES-C667-propietario}



\begin{quote}
Before the new one: say the Spanish for the torch, then the Spanish for the flat.
\end{quote}

\subsection*{You'll want to know: el propietario}
\textbf{el propietario} — \textquotedblleft{}the proprietor, the owner; the landlord\textquotedblright{}.

\begin{quote}
El propietario nos da las llaves.
\end{quote}

\begin{quote}
\emph{The owner gives us the keys.}
\end{quote}

A woman is \emph{la propietaria}. In speech \emph{el dueño} is common.

\begin{cousinweb}
\emph{propietario} is a learned borrowing from Latin \emph{proprietārius} \textquotedblleft{}an owner, a proprietor\textquotedblright{} (Lewis \& Short), built on \emph{proprius} \textquotedblleft{}own\textquotedblright{}.
\end{cousinweb}

\begin{grammarlens}[title={What you've built}]
One more word for telling, explaining and describing at B1.
\end{grammarlens}

\subsection*{Your turn}
\begin{itemize}
\raggedright
  \item \emph{Say it:} \emph{el propietario}
  \item \emph{Say it:} \emph{el propietario}, once more
  \item \emph{Say it:} a sentence of your own with \emph{el propietario}
  \item \emph{Recall:} say the Spanish for the torch, then the Spanish for the flat, then say \emph{el propietario} again
\end{itemize}

\begin{tcolorbox}[breakable,colback=teal!4,colframe=teal!35!black,title={Before you move on}]
What does el propietario mean? (\textquotedblleft{}The proprietor, the owner; the landlord\textquotedblright{}.) Say it once more.
\end{tcolorbox}
\section[el anfitrión]{el anfitrión — the host}

\label{lesson:ES-C667-anfitrion}



\begin{quote}
Before the new one: say the Spanish for the flat, then the Spanish for the proprietor.
\end{quote}

\subsection*{You'll want to know: el anfitrión}
\textbf{el anfitrión} — \textquotedblleft{}the host\textquotedblright{}.

\begin{quote}
El anfitrión nos espera en la puerta.
\end{quote}

\begin{quote}
\emph{The host is waiting for us at the door.}
\end{quote}

A woman is \emph{la anfitriona}.

\begin{cousinweb}
\emph{anfitrión} comes from the name \emph{Anfitrión}, Greek \emph{Amphitryōn} (EtymDB), a Theban hero of Greek legend; in an old comedy a great dinner is given at his house.
\end{cousinweb}

\begin{grammarlens}[title={What you've built}]
One more word for telling, explaining and describing at B1.
\end{grammarlens}

\subsection*{Your turn}
\begin{itemize}
\raggedright
  \item \emph{Say it:} \emph{el anfitrión}
  \item \emph{Say it:} \emph{el anfitrión}, once more
  \item \emph{Say it:} a sentence of your own with \emph{el anfitrión}
  \item \emph{Recall:} say the Spanish for the flat, then the Spanish for the proprietor, then say \emph{el anfitrión} again
\end{itemize}

\begin{tcolorbox}[breakable,colback=teal!4,colframe=teal!35!black,title={Before you move on}]
What does el anfitrión mean? (\textquotedblleft{}The host\textquotedblright{}.) Say it once more.
\end{tcolorbox}
\section[el inquilino]{el inquilino — the tenant}

\label{lesson:ES-C667-inquilino}



\begin{quote}
Before the new one: say the Spanish for the proprietor, then the Spanish for the host.
\end{quote}

\subsection*{You'll want to know: el inquilino}
\textbf{el inquilino} — \textquotedblleft{}the tenant\textquotedblright{}.

\begin{quote}
El inquilino paga el alquiler cada mes.
\end{quote}

\begin{quote}
\emph{The tenant pays the rent every month.}
\end{quote}

A woman is \emph{la inquilina}.

\begin{cousinweb}
\emph{inquilino} is a learned borrowing from Latin \emph{inquilīnus} \textquotedblleft{}a tenant, a lodger\textquotedblright{} (Lewis \& Short; EtymDB), from \emph{incolere} \textquotedblleft{}to dwell in\textquotedblright{}.
\end{cousinweb}

\begin{grammarlens}[title={What you've built}]
One more word for telling, explaining and describing at B1.
\end{grammarlens}

\subsection*{Your turn}
\begin{itemize}
\raggedright
  \item \emph{Say it:} \emph{el inquilino}
  \item \emph{Say it:} \emph{el inquilino}, once more
  \item \emph{Say it:} a sentence of your own with \emph{el inquilino}
  \item \emph{Recall:} say the Spanish for the proprietor, then the Spanish for the host, then say \emph{el inquilino} again
\end{itemize}

\begin{tcolorbox}[breakable,colback=teal!4,colframe=teal!35!black,title={Before you move on}]
What does el inquilino mean? (\textquotedblleft{}The tenant\textquotedblright{}.) Say it once more.
\end{tcolorbox}
\section[alquilar]{alquilar — to rent, to hire; to rent out}

\label{lesson:ES-C667-alquilar}



\begin{quote}
Before the new one: say the Spanish for the host, then the Spanish for the tenant.
\end{quote}

\subsection*{You'll want to know: alquilar}
\textbf{alquilar} — \textquotedblleft{}to rent, to hire; to rent out\textquotedblright{}.

\begin{quote}
Alquilamos un coche para tres días.
\end{quote}

\begin{quote}
\emph{We are hiring a car for three days.}
\end{quote}

\emph{Se alquila} on a sign — for rent.

\begin{cousinweb}
\emph{alquilar} is built on \emph{alquiler} \textquotedblleft{}rent\textquotedblright{} (EtymDB), which comes from Andalusian Arabic \emph{al-kirá} \textquotedblleft{}the hire, the rent\textquotedblright{}, with the Arabic article fused on.
\end{cousinweb}

\begin{grammarlens}[title={What you've built}]
One more word for telling, explaining and describing at B1.
\end{grammarlens}

\subsection*{Your turn}
\begin{itemize}
\raggedright
  \item \emph{Say it:} \emph{alquilar}
  \item \emph{Say it:} \emph{alquilar}, once more
  \item \emph{Say it:} a sentence of your own with \emph{alquilar}
  \item \emph{Recall:} say the Spanish for the host, then the Spanish for the tenant, then say \emph{alquilar} again
\end{itemize}

\begin{tcolorbox}[breakable,colback=teal!4,colframe=teal!35!black,title={Before you move on}]
What does alquilar mean? (\textquotedblleft{}To rent, to hire; to rent out\textquotedblright{}.) Say it once more.
\end{tcolorbox}
